\chapter*{三角比}
\subsection*{三姐妹}
\begin{enumerate}
    \item 已知 $ sin \alpha+ cos \alpha=\dfrac{1}{5}(0<\alpha<\pi)$, 求: $ sin \alpha- cos \alpha$
    \begin{jieda}
        $(sin\alpha+cos\alpha)^2 = 1 + 2sin\alpha cos\alpha = \dfrac{1}{25}$，因此$sin\alpha cos\alpha = -\dfrac{12}{25}$，根据$\alpha \in (0, \pi)$，所以$sin\alpha > 0$，因此$cos\alpha < 0$，即$\alpha$是第二象限角. \\
        因此有$(sin\alpha - cos\alpha)^2 = 1 - 2sin\alpha cos\alpha = \dfrac{49}{25}$ \\
        故$sin\alpha - cos\alpha = \pm \dfrac{7}{5}$, 根据$\alpha$是第二象限角，故$sin\alpha - cos\alpha = \dfrac{7}{5}$
    \end{jieda}
    \item (2025$\cdot$嘉定区二模$\cdot$6)已知 $\theta  \in  \mathbf{R}$ ，若 $\tan \theta  + \cot \theta  = 5$ ，则 $\sin {2\theta } =$ \blkda[2]{$\dfrac{2}{5}$}.
\end{enumerate}

\subsection*{诱导公式+给值求值类问题}
\begin{enumerate}
    \item (2024$\cdot$闵行区一模$\cdot$2)若 $\sin \alpha  = \dfrac{1}{3}$ ，则 $\sin \left( {\pi  - \alpha }\right)  =$ \blkda[2]{$\dfrac{1}{3}$}.

    \item (2024$\cdot$闵行区二模$\cdot$3)始边与 $x$ 轴的正半轴重合的角 $\alpha$ 的终边过点$(3,-4)$，则 $\sin \left( {\alpha  + \pi }\right)  =$ \blkda[2]{$\dfrac{4}{5}$}.
\end{enumerate}

\subsection*{用诱导公式化名}
\begin{enumerate}
    \item 已知$\alpha$的终边过点$(sin5, cos5)$，则$\alpha = $\blkda[3.5]{$\dfrac{\pi}{2} - 5 + 2k\pi, k \in \mathbf{Z}$}
\end{enumerate}


\subsection*{两角和与差+给值求值类问题}
\begin{enumerate}
    \item 已知锐角$\alpha, \beta$满足$cos\alpha = \dfrac{4}{5}, tan(\alpha - \beta) = -\dfrac{1}{3}$， 则$cos\beta=$ \blkda{$\dfrac{9\sqrt{10}}{50}$}
    \begin{jieda}
        首先判断出所求和已知的角之间的关系：$\beta = \alpha - (\alpha - \beta)$ \\
        $tan\beta = tan(\alpha - (\alpha - \beta)) = \dfrac{tan\alpha - tan(\alpha - \beta)}{1 + tan\alpha \cdot tan(\alpha - \beta)}$ \\
        $cos\alpha = \dfrac{4}{5} \Longrightarrow tan\alpha = \dfrac{3}{4}$，代入得到$tan\beta = \dfrac{\dfrac{3}{4} - \left(-\dfrac{1}{3}\right)}{1 + \dfrac{3}{4} \cdot \left (-\dfrac{1}{3} \right)} = \dfrac{13}{9}$\\
        故由$1 + tan^2\beta = \dfrac{1}{cos^2\beta}$，得到$cos\beta = \dfrac{9\sqrt{10}}{50}$ \\
        \ps \textbf{三角求值类问题，可以先从已知条件的角和未知的角之间的联系入手进行分析.}
    \end{jieda}
    \item (2024$\cdot$长宁区一模$\cdot$15)设点 $P$ 是以原点为圆心的单位圆上的动点，它从初始位置 ${P}_{0}\left( {1,0}\right)$ 出发，沿单位圆按逆时针方向转动角 $\alpha \left( {0 < \alpha  < \dfrac{\pi }{2}}\right)$ 后到达点 ${P}_{1}$ ，然后继续沿单位圆按逆时针方向转动角 $\dfrac{\pi }{4}$ 到达 ${P}_{2}$. 若点 ${P}_{2}$ 的横坐标为 $- \dfrac{3}{5}$ ，则点 ${P}_{1}$ 的纵坐标 \dotfill{(\kh{D})}
    \xx{$\dfrac{\sqrt{2}}{10}$}{$\dfrac{\sqrt{2}}{5}$}{$\dfrac{3\sqrt{2}}{5}$}{$\dfrac{7\sqrt{2}}{10}$}
    \begin{jieda}
        由题可知 $\cos \left( {\alpha  + \dfrac{\pi }{4}}\right)  =  - \dfrac{3}{5}$ ,且 $0 < \alpha  < \dfrac{\pi }{2}$ ,因为 $\dfrac{\pi }{4} < \alpha  + \dfrac{\pi }{4} < \dfrac{3\pi }{4}$ ,可知 $\dfrac{\pi }{2} < \alpha  + \dfrac{\pi }{4} < \dfrac{3\pi }{4}$ ,则 $\sin \left( {\alpha  + \dfrac{\pi }{4}}\right)  =  \sqrt{1 - {\cos }^{2}\left( {\alpha  + \dfrac{\pi }{4}}\right) } = \dfrac{4}{5}$ ,所以 $\sin \alpha  = \sin \left\lbrack  {\left( {\alpha  + \dfrac{\pi }{4}}\right)  - \dfrac{\pi }{4}}\right\rbrack   =  \sin \left( {\alpha  + \dfrac{\pi }{4}}\right) \cos \dfrac{\pi }{4} - \cos \left( {\alpha  + \dfrac{\pi }{4}}\right) \sin \dfrac{\pi }{4} = \dfrac{\sqrt{2}}{2}\left( {\dfrac{4}{5} + \dfrac{3}{5}}\right)  =  \dfrac{7\sqrt{2}}{10}$.
    \end{jieda}
\end{enumerate}

\subsection*{化幂与化角}
\begin{enumerate}
    \item 已知$\dfrac{3\pi}{2}<\theta<2\pi$, 化简: $\sqrt{\dfrac{1}{2}-\dfrac{1}{2}\sqrt{\dfrac{1}{2}+\dfrac{1}{2}\cos2\theta}}$ =\blkda{$\sin\dfrac{\theta}{2}$}.         
    \begin{jieda}
        $\sqrt{\dfrac{1}{2}+\dfrac{1}{2}cos2\alpha} = \sqrt{cos^2\alpha}$，因为$2\alpha \in (\dfrac{3\pi}{2}, 2\pi)$，所以$\sqrt{\dfrac{1}{2}+\dfrac{1}{2}cos2\alpha} = -cos\alpha$，\\
        $\sqrt{\dfrac{1}{2}-\dfrac{1}{2}cos\alpha} = \sqrt{sin^2\dfrac{\alpha}{2}}$，因为$2\alpha \in (\dfrac{3\pi}{2}, 2\pi)$，所以原式$=sin\dfrac{\theta}{2}$.
    \end{jieda}
    \item (2025$\cdot$金山区一模$\cdot$13)函数 $y = 1 - 2{\sin }^{2}x$ 是 \dotfill{(\kh{B})}
    \xx{最小正周期为 $\pi$ 的奇函数}{最小正周期为 $\pi$ 的偶函数}{最小正周期为 $\dfrac{\pi }{2}$ 的奇函数}{最小正周期为 $\dfrac{\pi }{2}$ 的偶函数}
\end{enumerate}

\subsection*{辅助角公式化简齐次三角函数解析式}
\begin{enumerate}
    \item (2025$\cdot$黄浦区二模$\cdot$6)函数 $f\left( x\right)  = \sqrt{3}\sin x + \sin \left( {\dfrac{\pi }{2} + x}\right)$ 的最大值是\blkda[2]{2}.
    \item (2025$\cdot$青浦区一模$\cdot$17)已知 $f\left( x\right)  = \left( {2{\cos }^{2}x - 1}\right) \sin {2x} + \dfrac{1}{2}\cos {4x}$.
    \begin{enumerate}
        \item 求函数 $y = f\left( x\right)$ 的最小正周期及最大值
        \item 若 $\alpha  \in  \left( {\dfrac{\pi }{2},\pi }\right)$ ，且 $f\left( \alpha \right)  = \dfrac{\sqrt{2}}{2}$ ，求 $\alpha$ 的值.
    \end{enumerate}
    \begin{jieda}
        \begin{enumerate}
            \item 因为 $f\left( x\right)  = \left( {2{\cos }^{2}x - 1}\right) \sin {2x} + \dfrac{1}{2}\cos {4x} =  \cos {2x}\sin {2x} + \dfrac{1}{2}\cos {4x} = \dfrac{1}{2}\left( {\sin {4x} + \cos {4x}}\right)  = \dfrac{\sqrt{2}}{2}\sin ({4x} +  \left. \dfrac{\pi }{4}\right)$ ,所以 $y = f\left( x\right)$ 的最小正周期为 $\dfrac{\pi }{2}$ ,最大值为 $\dfrac{\sqrt{2}}{2}$
            
            \item 因为 $\alpha  \in  \left( {\dfrac{\pi }{2},\pi }\right)$ ，所以 ${4\alpha } + \dfrac{\pi }{4} \in  \left( {\dfrac{9\pi }{4},\dfrac{17\pi }{4}}\right)$ . 因为 $f\left( \alpha \right)  =  \dfrac{\sqrt{2}}{2}$ ,所以 $f\left( \alpha \right)  = \dfrac{\sqrt{2}}{2}\sin \left( {{4\alpha } + \dfrac{\pi }{4}}\right)  = \dfrac{\sqrt{2}}{2}$ ,即 $\sin \left( {{4\alpha } + \dfrac{\pi }{4}}\right)  = 1$ . 所以 ${4\alpha } + \dfrac{\pi }{4} = \dfrac{5\pi }{2}$ ,故 $\alpha  = \dfrac{9\pi }{16}$ .
        \end{enumerate}
    \end{jieda}
\end{enumerate}

\subsection*{二倍角公式+给值求值类问题}
\begin{enumerate}
    \item 若$\dfrac{\sin\alpha+\cos\alpha}{\sin\alpha-\cos\alpha}=\dfrac{1}{2}$, 则$\tan2\alpha$=\blkda{$\dfrac{3}{4}$}. 
    \begin{jieda}
    齐次化后得$\tan\alpha=-3$, 故$\tan2\alpha=\dfrac{2\tan\alpha}{1-\tan^2\alpha}=\dfrac{3}{4}$. 
    \end{jieda}

    \item 若$sin(\dfrac{\pi}{6} - \alpha) = \dfrac{1}{3}$，则$cos(\dfrac{2\pi}{3} + 2\alpha) = $\blkda{$-\dfrac{7}{9}$}
        \begin{jieda}
            $cos(\dfrac{\pi}{3} + \alpha) = sin(\dfrac{\pi}{6} - \alpha) = \dfrac{1}{3}$，故$cos(\dfrac{2\pi}{3} + 2\alpha) = -\dfrac{7}{9}$
        \end{jieda}
\end{enumerate}

\subsection*{半角公式与解的个数问题}
\begin{enumerate}
    \item 若$sinx = \dfrac{4}{5}$，且$\dfrac{\pi}{2} < x < \pi$，则$sin\dfrac{x}{2} = $\blkda{$\dfrac{2\sqrt{5}}{5}$}
    \begin{jieda}
        $cosx = -\dfrac{3}{5}$，$\dfrac{x}{2} \in \left (\dfrac{\pi}{4}, \dfrac{\pi}{2} \right)$，故$sin\dfrac{x}{2} = \sqrt{\dfrac{1-cosx}{2}} = \sqrt{\dfrac{4}{5}} = \dfrac{2\sqrt{5}}{5}$
    \end{jieda}

    \item 已知$25sin^2\theta + sin\theta - 24 = 0$，$\theta$为第二象限角，则$cos\dfrac{\theta}{2} = $\blkda{$\dfrac{3}{5}$或$-\dfrac{3}{5}$}
    \begin{jieda}
        由$25sin^2\theta + sin\theta - 24 = 0$解得$sin\theta = -1$或$\dfrac{24}{25}$. \\
        由于$\theta$为第二象限角，所以$sin\theta = \dfrac{24}{25}, cos\theta = -\dfrac{7}{25}$，故$cos\dfrac{\theta}{2} = \pm\sqrt{\dfrac{1+cos\theta}{2}} = \pm\dfrac{3}{5}$
    \end{jieda}
\end{enumerate}

\subsection*{半角公式与万能公式}
\begin{enumerate}
    \item 若方程$x^2+4ax+3a+1 = 0 (a>1)$的两根分别是$tan\alpha, tan\beta$，且$\alpha, \beta \in (-\dfrac{\pi}{2}, \dfrac{\pi}{2})$，求$tan(\dfrac{\alpha+\beta}{2})$.
    \begin{jieda}
        根据韦达定理有$\left \{ \begin{array}{l}
                tan\alpha + tan\beta = -4a  \\
                tan\alpha \cdot tan\beta = 3a+1 
        \end{array}\right.$，因为$a > 1$，所以$tan\alpha, tan\beta < 0$，因此$\alpha, \beta \in (-\dfrac{\pi}{2}, 0)$ \\
        根据和角公式有$tan(\alpha + \beta) = \dfrac{tan\alpha + tan\beta}{1 - tan\alpha \cdot tan\beta} = \dfrac{4}{3}$，故$cos(\alpha+\beta) = -\dfrac{3}{5}, sin(\alpha + \beta) = -\dfrac{4}{5}$，故$tan(\dfrac{\alpha+\beta}{2}) = \dfrac{sin(\alpha+\beta)}{1+cos(\alpha + \beta)} = -2$. \\
        \ps 判别式大于零影响$a$的范围，但是对于本题$tan(\alpha + \beta)$的值不含$a$，所以对解题没影响.
    \end{jieda}

    \item 若$\dfrac{2sin\theta + cos\theta}{sin\theta - 3cos\theta} = -5$，则$3cos2\theta + 4sin2\theta = $\blkda{$\dfrac{7}{5}$}
    \begin{jieda}
        $\dfrac{2sin\theta + cos\theta}{sin\theta - 3cos\theta} = -5$即$\dfrac{2tan\theta + 1}{tan\theta - 3} = -5$，解得$tan\theta = 2$ \\
        故$cos2\theta = \dfrac{1-tan^2\theta}{1+tan^2\theta} = \dfrac{-3}{5}, sin2\theta = \dfrac{2tan\theta}{1+tan^2\theta} = \dfrac{4}{5}$，故$3cos2\theta + 4sin2\theta = \dfrac{7}{5}$
    \end{jieda}
\end{enumerate}

\subsection*{辅助角公式与$y = \dfrac{asinx+bcosx+c}{a'sinx+b'cosx+c'}$的值域}
\begin{enumerate}
    \item 求$y = \dfrac{sinx}{cosx - 3}$的值域.
    \begin{jieda}
        \begin{enumerate}
            \item 辅助角公式 \\
            去分母：$y (cosx - 3) = sinx$，即$sinx - y cosx = -3y$ \\
            应用辅助角公式：$\sqrt{1+y^2}sin(x+\varphi) = -3y$,$cos\varphi = \dfrac{1}{\sqrt{1+y^2}}, sin\varphi = \dfrac{-y}{\sqrt{1+y^2}}$ \\
            解不等式：$\sqrt{1+y^2} \geq |-3y|$，即$1+y^2 \geq 9y^2$，故$y \in \left[-\dfrac{\sqrt{2}}{4}, \dfrac{\sqrt{2}}{4}\right]$ \\
            \ps 定义域不是全体实数不建议使用，定义域不是自然定义域一定不能使用
            \item 万能公式
            当$x = k\pi (k \in \mathbf{Z})$时，$y = 0$
            
            当$x \neq k\pi (k \in \mathbf{Z})$时，$y = \dfrac{\dfrac{2tan\dfrac{x}{2}}{1+tan^2\dfrac{x}{2}}}{\dfrac{1-tan^2\dfrac{x}{2}}{1+tan^2\dfrac{x}{2}} - 3} = \dfrac{2tan\dfrac{x}{2}}{-2-4tan^2\dfrac{x}{2}} = -\dfrac{tan\dfrac{x}{2}}{1+2tan^2\dfrac{x}{2}}$，令$t = tan\dfrac{x}{2}, t \in \mathbf{R}$，故$y = -\dfrac{t}{1+2t^2}$
            \begin{dingautolist}{172}
                \item 当$t = 0$时，$y = 0$
                \item 当$t \neq 0$时，$y = - \dfrac{1}{2t+\dfrac{1}{t}} \in \left[-\dfrac{\sqrt{2}}{4}, 0\right) \cup \left(0,\dfrac{\sqrt{2}}{4}\right]$
            \end{dingautolist}
            综上，$y \in \left[-\dfrac{\sqrt{2}}{4}, \dfrac{\sqrt{2}}{4}\right]$
            \item 数形结合 \\
            $y$即点$(cosx, sinx)$到$(3, 0)$连线斜率
        \end{enumerate}
    \end{jieda}
\end{enumerate}